\section{给参数也建立概率模型}

似然只比较固定数据在不同参数下的相对支持。若要问“观察后参数落在某区间的概率”，必须再给参数一个概率分布，这正是 Bayes 模型增加的假设。

\begin{definition}[先验、后验与可信集]
设参数 \(\Theta\) 是随机变量，先验密度为 \(\pi(\theta)\)，给定 \(\Theta=\theta\) 时数据的密度为 \(p_\theta(x)\)。若
\(0<m(x)=\int p_\vartheta(x)\pi(\vartheta)\,d\vartheta<\infty\)，则观测 \(x\) 后的后验密度为
\[
\pi(\theta\mid x)=\frac{p_\theta(x)\pi(\theta)}{m(x)}.
\]
后验可信集 \(C(x)\) 满足
\(\int_{C(x)}\pi(\theta\mid x)\,d\theta=1-\alpha\)。
\end{definition}

分母是使后验归一化的边际数据密度。先验改变，后验一般也改变；若只写“后验与似然成正比”而不指定先验，推断尚未确定。可信集对随机参数给条件概率；频率学派置信集则要求对每个固定 \(\theta\)，重复抽样的覆盖概率达到指定值，两句话的随机对象不同。

取独立的 \(X_i\mid p\sim\operatorname{Bernoulli}(p)\)，先验
\(p\sim\operatorname{Beta}(a,b)\)，其中 \(a,b>0\)，密度正比于
\(p^{a-1}(1-p)^{b-1}\)。
记 \(s=\sum_iX_i\)，把似然 \(p^s(1-p)^{n-s}\) 乘上先验，逐项合并幂次，得
\[
p\mid X\sim\operatorname{Beta}(a+s,b+n-s).
\]
这里 \(\operatorname{Beta}(a,b)\) 不是只由比例号定义的记号。其密度为
\(p^{a-1}(1-p)^{b-1}/B(a,b)\)，\(0<p<1\)，其中
\(B(a,b)=\int_0^1u^{a-1}(1-u)^{b-1}\,\dd u\)；当
\(a,b>0\) 时两端的幂都可积。把二项似然乘上这个归一化密度，再对 \(p\) 积分，数据的边际概率是
\[
P(X=x)=\frac{B(a+s,b+n-s)}{B(a,b)}
\]
（这里 \(x\) 表示一条指定的零一序列；若只观测总成功数 \(s\)，右侧还需乘组合数 \(\binom ns\)）。因此后验归一化常数也被明确算出。分母的 \(B(a,b)\) 不依赖数据，而分子的两个参数分别累加了成功与失败的次数。

Beta 均值可直接从定义算，不必背表。分部积分应用于
\(u^a(1-u)^b\) 的导数，端点项为零，给
\(aB(a,b+1)=bB(a+1,b)\)；又由
\(B(a,b)=B(a+1,b)+B(a,b+1)\)，消元得
\(B(a+1,b)/B(a,b)=a/(a+b)\)。把 \(a,b\) 换成后验参数，便得到下面的后验均值。先验参数 \(a,b\) 可以直观地读成给成功与失败各加的权重，但它们不必是整数，也不等于实际额外做了 \(a+b\) 次试验。
在平方损失下，使后验期望损失
\(\mathbb E[(d-p)^2\mid X]\) 最小的 \(d\) 是后验均值：展开平方或对 \(d\) 求导均得
\(d=(a+s)/(a+b+n)\)。这说明小样本时先验权重不会自动消失。
例如只观察一次成功，均匀先验给后验 \(\operatorname{Beta}(2,1)\)，均值 \(2/3\)；而似然 \(L(p)=p\) 本身不是这个后验密度。

\begin{proposition}[后验预测]
上述 Beta--Bernoulli 模型中，若下一次试验 \(X_{n+1}\) 在给定 \(p\) 后与旧数据独立，则
\(P(X_{n+1}=1\mid X)=\mathbb E[p\mid X]=(a+s)/(a+b+n)\)。
\end{proposition}
\begin{proof}
对后验积分：
\(P(X_{n+1}=1\mid X)=
\int P(X_{n+1}=1\mid p,X)\pi(p\mid X)\,dp
=\int p\pi(p\mid X)\,dp\)。
第二个等号正用到给定 \(p\) 后的条件独立。
\end{proof}

有限样本公式是精确的。后验近似正态的 Bernstein--von Mises 结论属于另一层问题：它需要有限维可识别模型、真参数位于内部、正的 Fisher 信息、局部二次似然展开和先验在真值附近正且连续等条件。边界、混合模型或不可识别模型不能直接套用。这里不把渐近口号当成上面有限样本结论的证明。

\begin{exercise}
比较 \(a=b=1\) 与 \(a=b=100\) 时，\(n=5,s=4\) 的后验均值。说明两者为何不应被解释为“相同数据给出相同参数概率”。
\end{exercise}
